% Round 22 / geometry branch. Candidate; independent review required.
\section{Two curve coefficients and exact second monodromy}
\label{r22dw:section}

\paragraph{Scope.}
Keep the actual local datum and model of DS22.  The positivity
hypothesis on $S\times T$ is not needed in this section.  We compute
both entire extreme terms of the weight spectral sequence on $H^4$
and the actual map $N^2$ between them.  The computation retains the
curve blowups of the common double and their nonzero exceptional
restriction classes.  These corrections happen to vanish on a
specified tensor coefficient; they are not discarded formally.

All cohomology is geometric $\ell$-adic cohomology for
$\ell\ne\operatorname{char}k$.  Put
\[
\begin{aligned}
 V_S&=H^1(S),&W_C&=H^1(C),&
 P_C&=\ker\{i_{C*}:W_C\longrightarrow H^3(S)(1)\},\\
 V_T&=H^1(T),&W_D&=H^1(D),&
 P_D&=\ker\{i_{D*}:W_D\longrightarrow H^3(T)(1)\}.
\end{aligned}
 \tag{DW22.1}\label{r22dw:coefficients}
\]
The two cohomological hard Lefschetz isomorphisms
$i_{C*}i_C^*:V_S\to H^3(S)(1)$ and
$i_{D*}i_D^*:V_T\to H^3(T)(1)$ give Frobenius-equivariant splittings
\[
 \begin{aligned}
 W_C&=i_C^*V_S\oplus P_C,&
 W_D&=i_D^*V_T\oplus P_D,\\
 p_C&:=\dim P_C=2g(C)-b_1(S),&
 p_D&:=\dim P_D=2g(D)-b_1(T).
 \end{aligned}
 \tag{DW22.2}\label{r22dw:split}
\]
Indeed $i_C^*(i_{C*}i_C^*)^{-1}i_{C*}$ is the displayed projector,
and the second curve is identical.  No semisimplicity assertion or
algebraic inverse to these cohomological maps is needed.  This is
the same elementary splitting used in WD21, now in two independent
curve factors.

\subsection{The actual blowup corrections on the common double}

Write $W=C\times D\times\mathbf P^1$,
$E_j=C\times D\times\{p_j\}$, and
$Z_j=E_j\cap V(\bar Q)$.  Let
\[
 q:\widetilde W=\operatorname{Bl}_{\coprod_j Z_j}W\longrightarrow W,
 \quad j_j:F_j\hookrightarrow\widetilde W,
 \quad p_j:F_j\longrightarrow Z_j.
 \tag{DW22.3}\label{r22dw:blowup}
\]
The strict transform $E'_j$ is naturally isomorphic to $E_j$,
because $Z_j$ is Cartier in $E_j$.  Let $i_j:E_j\hookrightarrow W$
and $i'_j:E'_j\hookrightarrow\widetilde W$ denote the inclusions.
The common double of DS22 is this $\widetilde W$, and its adjacent
triples are precisely the $E'_j$.

\begin{lemma}[Actual restriction and Gysin identities]
\label{r22dw:actualmaps}
For $a\in\operatorname{CH}^1(W)_{\mathbf Q}$,
$z\in\operatorname{CH}^1(E_j)_{\mathbf Q}$, and scalars $\lambda_k$,
one has
\[
\begin{aligned}
 (i'_j)^*\left(q^*a+\sum_k\lambda_k[F_k]\right)
       &=i_j^*a+\lambda_j[Z_j],\\
 (i'_j)_*z&=q^*(i_j)_*z-(j_j)_*p_j^*(z|_{Z_j}).
\end{aligned}
 \tag{DW22.4}\label{r22dw:Chowmaps}
\]
The same formulas hold for all cohomological classes in the
corresponding degrees, including nonalgebraic classes.  Moreover
\[
 [Z_j]=\operatorname{pr}_C^*c_1(\mathcal A|_C)
                   +\operatorname{pr}_D^*c_1(\mathcal B|_D)
                         \quad\text{on }E_j=C\times D.
 \tag{DW22.5}\label{r22dw:external}
\]
In particular the exceptional restriction term has zero projection
to $P_C\otimes P_D\subset H^2(C\times D)$, and
$z|_{Z_j}=0$ in $H^2(Z_j)$ for every
$z\in W_C\otimes W_D$.
\end{lemma}

\begin{proof}
Since the centers are disjoint, it suffices to work with one.  The
strict transform of a Cartier divisor containing the center has
total transform $q^*E_j=E'_j+F_j$.  The Rees charts of the ideal
$(E_j,\bar Q)$ show that $F_j\cap E'_j$ is the section $Z_j$ of
the projective bundle $F_j\to Z_j$, and that the intersection is
transverse.  Thus $[F_j]|_{E'_j}=+[Z_j]$; the sign is positive.
The other exceptional divisors miss $E'_j$.  This proves the first
formula at the actual Chow level.

For the second formula, the two classes on its left and first
right term have the same proper pushforward to $W$.  The actual
codimension-two Chow blowup decomposition therefore writes their
difference uniquely as $(j_j)_*p_j^*b$ with
$b\in\operatorname{CH}^1(Z_j)_{\mathbf Q}$.  There are no terms
from other centers, since those exceptional divisors are disjoint
from $E'_j$.  Restrict to $F_j$.  The restriction of
$q^*(i_j)_*z$ is zero: it factors through a codimension-two class
on the curve $Z_j$.  The other restriction is
$s_*(z|_{Z_j})$, where $s:Z_j\hookrightarrow F_j$ is the indicated
section.  If $\xi=c_1(\mathcal O_{F_j}(1))$, then
$s_*\beta=\xi p_j^*\beta$ for $\beta\in\operatorname{CH}^1(Z_j)$:
the possible additional base-divisor term has product zero on a
curve.  Finally $j_j^*j_{j*}=-\xi$ and $p_{j*}\xi=1$.
Applying $p_{j*}$ identifies $b=-z|_{Z_j}$ and gives the second
formula.  This uses the proved direct blowup decomposition and
projective-bundle identity.  It does not cancel an arbitrary
Gysin map or assign a native degree to absolute total-model
$\operatorname{CH}_0$.

The cohomological blowup decomposition, transverse base change,
and the same projective-bundle identities give precisely this
proof for all geometric cohomology classes.  Equation
\eqref{r22dw:external} follows from the actual restricted line
bundle of $Q$: the $\mathcal O(d)$ factor becomes a constant
one-dimensional vector space on the native point $p_j$.
It is an equality of divisor classes even though the curve $Z_j$
itself need not be a product or a union of external divisors.

For $z\in W_C\otimes W_D$, cup product with either external
summand of \eqref{r22dw:external} is zero by the dimensions of the
curves.  Hence $\int_{Z_j}z|_{Z_j}=0$.  The curve $Z_j$ is
geometrically integral by hypothesis, and its $H^2$ is
one-dimensional with nonzero trace.  It follows that $z|_{Z_j}=0$.
The projection assertion is immediate from the Kunneth
decomposition of $H^2(C\times D)$.
\end{proof}

\subsection{Both complete extreme terms}

Write $D^{[m]}$ for the disjoint union of $m$-fold strata.  We use
Saito's convention, translated from zero-based to this one-based
stratum index:
\[
 E_1^{a,b}=\bigoplus_{i\ge\max(0,-a)}
 H^{b-2i}\bigl(D^{[a+2i+1]},\mathbf Q_\ell\bigr)(-i)
                 \Longrightarrow H^{a+b}(X_{\bar K},\mathbf Q_\ell).
 \tag{DW22.6}\label{r22dw:RZ}
\]
The $d_1$ maps are the signed actual restrictions and Gysin maps.
Because there are no quadruples, the complete terms affecting the
lower extreme are
\[
 E_1^{1,2}=H^2(D^{[2]}),\qquad
 E_1^{2,2}=H^2(D^{[3]})=H^2(C\times D)^{\oplus d},
 \qquad E_1^{3,2}=0.
 \tag{DW22.7}\label{r22dw:lowerterms}
\]
There are no omitted $H^0$-summands from fourfold intersections.
Likewise the full upper terms are
\[
 E_1^{-3,6}=0,\qquad
 E_1^{-2,6}=H^2(C\times D)^{\oplus d}(-2),\qquad
 E_1^{-1,6}=H^4(D^{[2]})(-1).
 \tag{DW22.8}\label{r22dw:upperterms}
\]

\begin{proposition}[Exact coefficient quotients and kernels]
\label{r22dw:extremes}
Put $P=P_C\otimes P_D$.  There are Frobenius-equivariant
identifications of the whole displayed terms
\[
\begin{aligned}
 E_2^{2,2}&=P\otimes
     \bigl(\mathbf Q_\ell^d/\mathbf Q_\ell(1,\ldots,1)\bigr),\\
 E_2^{-2,6}&=
   \left(P\otimes\ker\{\mathbf Q_\ell^d
                         \xrightarrow{\ \sum\ }\mathbf Q_\ell\}\right)(-2).
\end{aligned}
 \tag{DW22.9}\label{r22dw:extremeformula}
\]
In particular both have dimension $p_Cp_D(d-1)$.
\end{proposition}

\begin{proof}
For the lower extreme, a triple labeled $j$ has two private
doubles $C\times T$ and $S\times D$.  Their restriction images
in $H^2(C\times D)$ contain both external Kunneth terms
$H^2(C)\otimes H^0(D)$ and $H^0(C)\otimes H^2(D)$: restriction
from a surface to its ample curve is nonzero on an ample divisor
class.  Their mixed images are respectively
$W_C\otimes i_D^*V_T$ and $i_C^*V_S\otimes W_D$.
The splittings \eqref{r22dw:split} therefore identify the quotient
by these two private images with exactly $P$, independently for
each $j$.

The remaining common double is $\widetilde W$.
Its full $H^2$ is $q^*H^2(W)\oplus\bigoplus_j\mathbf Q_\ell[F_j]$.
The pullback summand maps to the same $H^2(C\times D)$ class in
all $d$ triples; the $\mathbf P^1$ class restricts to zero.  By
Lemma~\ref{r22dw:actualmaps}, the exceptional class $[F_j]$
restricts to $[Z_j]$ in its own coordinate and to zero elsewhere.
These external divisor classes already lie in the private images.
Consequently the common double contributes precisely the diagonal
$P\to P^d$ to the remaining quotient.  With the ordering
$0<1<j+1$, its incidence sign is always the positive sign of
the face $\{0,1\}$ in $\{0,1,j+1\}$.  This proves the first
formula and shows it is the whole lower term.

For the upper extreme, consider a coordinate
$z_j\in H^2(C\times D)$.  Its Gysin map to $C\times T$ sends
the two external terms into, respectively,
$H^2(C)\otimes H^2(T)$ and $H^0(C)\otimes H^4(T)$.
These maps are injective on their one-dimensional domains:
$[D]\ne0$ and the point trace is nonzero.  The mixed term goes
by $1\otimes i_{D*}$ into $W_C\otimes H^3(T)$.
The Kunneth summands are distinct, so the kernel of this private
Gysin is exactly $W_C\otimes P_D$.  The other private Gysin
has kernel exactly $P_C\otimes W_D$.  Each private double has
only this one adjacent triple.  Thus their two conditions force
each $z_j$ individually into $P$, with no cancellation across
different triples.

For $z_j\in P$, the exceptional term in
\eqref{r22dw:Chowmaps} is zero in cohomology by the same lemma.
Its Gysin to the common double is therefore
$q^*(z_j\boxtimes h)$, where $h$ is the degree-one point class
of $\mathbf P^1$.  The map $z\mapsto q^*(z\boxtimes h)$ is
injective by Kunneth and the blowup decomposition.  Its incidence
sign is common to all $j$, since the omitted point vertex is
always in the same ordered position.  The remaining kernel
condition is consequently $\sum_j z_j=0$.  There is no incoming
quotient in \eqref{r22dw:upperterms}.  This proves the second
formula and its asserted completeness.
\end{proof}

\subsection{The actual map and a sharp dominance-depth obstruction}

\begin{theorem}[Exact second monodromy, DW22]
\label{r22dw:monodromy}
On the actual generic cohomology one has
\[
 \operatorname{rank}\!\left(N^2:H^4(X_{\bar K},\mathbf Q_\ell)
                    \longrightarrow H^4(X_{\bar K},\mathbf Q_\ell)(-2)\right)
                       =p_Cp_D(d-1),\qquad N^3=0.
 \tag{DW22.10}\label{r22dw:rank}
\]
More precisely,
$N^2:E_2^{-2,6}\to E_2^{2,2}(-2)$ is an isomorphism.  These
statements survive every finite valued extension.  If
$p_Cp_D(d-1)>0$, the minimum crossing depth, taken over all finite
valued extensions and all projective reduced strictly semistable
witnesses with smooth projective generic fiber dominating $X$, is
exactly two.  The lower bound holds over every fixed finite valued
extension; the original model supplies the upper bound for this
minimum.  Here depth means the maximum number of components through
a point minus one.
\end{theorem}

\begin{proof}
Saito, Corollary 2.2.4(2), identifies the degree-two shift of a
tame inertia operator $T-1$ on \eqref{r22dw:RZ} with
$t_\ell(T)$ times the identity on the common stratum summands.
The normalized logarithm $N=\log(T)/t_\ell(T)$ has the same
graded map: higher terms of the logarithm lower the filtration
by at least four.  Applying it twice sends the common triple
$H^2$ summand in column $-2$ to that in column $2$ by the
identity, with the displayed Tate twist.  Under
\eqref{r22dw:extremeformula}, its induced map is therefore the
tensor product with $P$ of
\[
 \ker(\mathbf Q_\ell^d\xrightarrow{\sum}\mathbf Q_\ell)
       \longrightarrow\mathbf Q_\ell^d/\mathbf Q_\ell(1,\ldots,1).
 \tag{DW22.11}\label{r22dw:scalariso}
\]
This map is an isomorphism: a constant vector in the kernel has
sum $dc=0$, hence $c=0$, and both dimensions equal $d-1$.
Equivalently it comes from the rational direct sum into the
diagonal line and the sum-zero hyperplane.  There is no use of
an ordering or positivity on a $\mathbf Q_\ell$-valued form.

The spectral sequence degenerates at $E_2$: each $E_1^{a,b}$
is pure of weight $b$ by finite-field purity, so each $E_2^{a,b}$
has that same weight; all higher differentials have distinct
source and target weights.  The abutment filtration on $H^4$
has weights in $[2,6]$, since only columns $-2\le a\le2$
occur.  The operator $N$ lowers this filtration by two.
Thus $N^2$ kills $W_5$, lands in $W_2(-2)$, and is completely
the just-computed map between the extreme graded pieces,
because $W_1=0$.  This proves the exact rank, not merely a
lower bound.  A third application lands in $W_0=0$, proving
$N^3=0$ without a general weight--monodromy assumption.

Restriction to a finite valued extension multiplies the logarithmic
monodromy by its positive ramification index when the two tame
characters are normalized using their respective uniformizers.
With the transported character the operator is unchanged instead;
only its invariance up to a nonzero scalar is used here.  All ranks
and nilpotence assertions therefore persist.  Suppose a smooth
projective $Y$ dominates $X$ over such an extension and has a
reduced strictly semistable model of depth at most one.  Its weight
spectral sequence has only columns $-1,0,1$, so $N^2=0$ on
$H^4(Y)$.  Pullback $H^4(X)\to H^4(Y)$ is injective and commutes
with inertia.  Indeed, if $m=\dim Y\ge4$ and $h$ is an ample
class on $Y$, then proper trace followed by cup product with
$h^{m-4}$ satisfies
\[
 f_*(h^{m-4}f^*x)=c\,x,
 \qquad c=\deg(h^{m-4}|_{Y_{\eta_X}})>0.
 \tag{DW22.12}\label{r22dw:dominanttrace}
\]
For a generically finite morphism $c$ is its positive degree.
This is an ordinary cohomological projection formula; no actual
analytic proper pushforward is inferred.  Injectivity contradicts
the nonzero $N^2$ of $X$.  The actual DS22 model itself has depth
two over $K$, and the identity is a dominating witness.  Hence the
minimum over all such fields and witnesses is precisely two.
No semistable upper-bound assertion over each fixed ramified
extension is needed for this conclusion.
\end{proof}

\paragraph{A coefficient-sensitive nonarchimedean distinction.}
The ordinary dual complex of this actual model is the edge
$\{0,1\}$ joined with $d$ isolated point vertices.  It is a union
of $d$ triangles glued along one common edge, hence is contractible.
Nevertheless \eqref{r22dw:rank} is nonzero whenever $p_C,p_D>0$.
The nonzero information lies in the actual coefficient
$P_C\otimes P_D$ and in its restriction/Gysin complex, not in
constant-coefficient cohomology of that dual complex.  DS22 proves
full G for the explicit positive components even in this situation.
Thus neither this coefficient group nor nonzero $N^2$ is by itself
a Green obstruction.  What a general theory still must control is
the cycle-visible period map, not annihilate every nearby-cycle
coefficient.

\paragraph{Primary-source check and dependencies.}
The new work is the exact actual blowup correction and the complete
two-curve coefficient calculation.  The spectral sequence and
its differential are Saito, \emph{Weight spectral sequences and
independence of $\ell$}, June 9, 2003, Corollary 2.2.4 and
Proposition 2.2.6, at
\url{https://www.ms.u-tokyo.ac.jp/~t-saito/pp/wmr2r.pdf}.
The pure stratum weights are Deligne, \emph{La conjecture de Weil I},
Theorem 1.6, at
\url{https://www.numdam.org/item/PMIHES_1974__43__273_0.pdf}.
The curve splittings use cohomological hard Lefschetz as in
Deligne, \emph{La conjecture de Weil II}, Theorem 4.1.1 and
Lemma 4.1.2, at
\url{https://publications.ias.edu/sites/default/files/Number40.pdf}.
The independent frozen WD21 source uses the same monodromy
normalization.  No new literature-priority claim is made.
